\documentclass[UTF8]{ctexart}

\usepackage{amsmath, amssymb, amsthm}
\usepackage{mathrsfs}
\usepackage{ctex}
\usepackage{tikz}
\usepackage{tikz-cd}
\usepackage{pgfplots}
\usepackage{xltabular}
\usepackage{booktabs}
\usepackage{geometry}
\usepackage{extarrows}
\geometry{a4paper, margin=2cm}
\pgfplotsset{compat=1.18}
\usetikzlibrary{shapes.geometric, positioning, calc, decorations.pathreplacing,3d ,positioning}

\newtheorem{theorem}{定理}[subsection]
\newtheorem{proposition}[theorem]{命题}
\newtheorem{corollary}[theorem]{推论}
\newtheorem{lemma}[theorem]{引理}
\theoremstyle{definition}
\newtheorem{definition}[theorem]{定义}
\newtheoremstyle{examplestyle}{3pt}{3pt}{\itshape}{0pt}{\bfseries}{.}{.5em}{\thmnumber{例 \arabic{example}}\thmnote{\ (#3)}}
\theoremstyle{examplestyle}
\newtheorem{example}{}[subsection]
\theoremstyle{remark}
\newtheorem*{remark}{注}

\title{L1凸松弛压缩感知信号恢复}
\date{}

\begin{document}
研究问题: 带噪声稀疏信号恢复
\[y=Ax_{\text{true}}+e\]
\begin{table}[htbp]
\centering
\begin{tabular}{ll}
\hline
符号 & 含义 \\
\hline
$x_{\text{true}}\in\mathbb{R}^n$ & $K$-稀疏原始信号 \\
$A\in\mathbb{R}^{M\times n}$ & 高斯随机测量矩阵 \\
$e$ & 高斯观测噪声 \\
$y$ & 观测向量 \\
\hline
\end{tabular}
\caption{符号说明}
\label{tab:notation}
\end{table}

原始问题:L0稀疏恢复
\[\begin{aligned}
    &\text{min}\quad &&\|x\|_0\\
    &\text{s.t.}\quad &&Ax=y
\end{aligned}\]
但是零范数非凸, 不连续, 难以大规模精确求解, 凸松弛为一范数. 同时现实观测存在噪声, 等式约束无法严格满足, 替换为残差约束.
\[\begin{aligned}
    &\text{min}\quad &&\|x\|_1\\
    &\text{s.t.}\quad &&\|Ax-y\|_2\le\delta
\end{aligned}\]
$\delta$取值影响优化问题可行域, 取太小可行域空, 取太大则约束过于松弛, 得到的解过度稀疏, 失真严重.

对于观测模型
\[y=Ax_{\text{true}}+e,\quad e\sim N(0,\sigma^2I_M)\]
残差向量$r=Ax-y=A(x-x_{\text{true}})-e$, 当$x=x_\text{true}$时, 真实残差满足
\[r_{\text{true}}=-e\sim N(0,\sigma^2I_M)\]
标准化得
\[\frac{1}{\sigma}r_{\text{true}}\sim N(0,I_M)\]
于是
\[Z=\frac{1}{\sigma^2}\|r_\text{true}\|_2^2=\frac{1}{\sigma^2}\|e\|_2^2\sim\chi^2(M)\]
我们可以取卡方分布的高概率上界$\sigma\sqrt{M+2\sqrt{2M}}$为$\delta$.

现在重新审视目标问题, 目标函数$\|x\|_1$是凸函数, 约束函数$\|Ax-y\|_2$是凸函数. 所以这是一个凸优化问题, 局部最优解等价于全局最优解. 下面推导KKT条件.

Lagrange函数:
\[L(x,\lambda)=\|x\|_1+\lambda(\|Ax-y\|_2-\delta)\]
由于$\|x\|_1$在原点不可微, 稳定性条件中梯度用次梯度代替. $\partial \|x\|_1=s$, 其中$s$满足
\[s_1\in\text{Sign}(x_i)=
\begin{cases}
    \{1\},\quad &x_i>0\\
    \{-1\},\quad &x_i<0\\
    [-1,1],\quad &x_i=0
\end{cases}\]
而二范数可微, 次梯度就是梯度
\[\partial \|Ax-y\|_2=\frac{A^T(Ax-y)}{\|Ax-y\|_2}\]

所以KKT条件如下

原问题可行:$\|Ax-y\|_2\le\delta$.

对偶问题可行:$\lambda\ge0$.

互补松弛:$\lambda(\|Ax-y\|_2-\delta)=0$.

稳定性条件:$-\lambda\frac{A^T(Ax-y)}{\|Ax-y\|_2}\in\partial\|x\|_1$

其中互补松弛项说明, 要么解在约束边界$\|Ax-y\|_2=\delta$上($\lambda>0$时), 要么$\lambda=0$, 约束无效, 等价无约束最小化$\|x\|_1$.

对偶问题：
\[\begin{aligned}
    &\text{min} &&y^Tv+\delta\|v\|_2\\
    &\text{s.t.} &&\|A^Tv\|_\infty\le1
\end{aligned}\]
强对偶成立.
\end{document}